% ============================================================
% Come si usano, argomento 5: variabili aleatorie discrete
% ============================================================
\argomento{Variabili aleatorie discrete}{Formule a p.~\pageref*{f:pmf}}

\begin{scelta}{Come scegliere la variabile}
\renewcommand{\arraystretch}{1.25}
\begin{tabular}{@{}p{0.6\linewidth}p{0.34\linewidth}@{}}
  \toprule
  \textbf{Se si conta o si osserva\ldots} & \textbf{la variabile è\ldots} \\
  \midrule
  l'esito di una sola prova con due possibilità (bit giusto o sbagliato, testa o croce)
    & di Bernoulli \vv{bernoulli} \\
  il numero di successi in $N$ prove indipendenti, tutte con la stessa probabilità $p$
    & binomiale \vv{distbinomiale} \\
  un valore fra $M$, tutti ugualmente probabili (dado onesto, byte casuale)
    & uniforme \vv{uniforme} \\
  il numero di eventi rari in un intervallo, con media $\lambda$ nota (arrivi, guasti, errori); oppure una binomiale
    con $N$ grande e $p$ piccolo
    & di Poisson \vv{poisson}, con $\lambda = N p$ \\
  \bottomrule
\end{tabular}

Per qualsiasi variabile lo schema è lo stesso: si scrive l'alfabeto, si scrive la pmf \vv{pmf}, si controlla che sommi
a 1, si calcola la media \vv{media}.
\end{scelta}

\scheda{pmf}{
  \usoquando{Per descrivere una variabile aleatoria discreta: la pmf dice con che probabilità assume ciascun
    valore.}
  \usocome{Per trovare la pmf di una variabile definita sui risultati di un esperimento, per ogni valore $x$ si
    sommano le probabilità dei risultati che la variabile manda in $x$. La probabilità di un insieme di valori è la
    somma dei valori della pmf.}
  \usoattenzione{Una sequenza è una pmf solo se tutti i valori sono $\geq 0$ \emph{e} sommano a 1:
    $(0{,}6;\ 0{,}6;\ -0{,}2;\ 0)$ somma a 1, ma non è una pmf.}
  \usoesempio{Due lanci di una moneta, $Y$ = numero di croci: $p_Y(0) = P(\{TT\}) = \frac{1}{4}$,
    $p_Y(1) = P(\{TC, CT\}) = \frac{1}{2}$, $p_Y(2) = P(\{CC\}) = \frac{1}{4}$.}
}

\scheda{mediacamp}{
  \usoquando{Si hanno $n$ valori osservati (dati, simulazioni) e se ne vuole il valore medio.}
  \usocome{Si sommano e si divide per $n$; con molti valori ripetuti conviene raggruppare, moltiplicando ogni
    valore per la sua frequenza relativa.}
  \usoesempio{Venti lanci di un dado con somma 75: $\overline{X}_{20} = 75/20 = 3{,}75$. Al crescere dei lanci si
    avvicina alla media statistica, $3{,}5$.}
}

\scheda{media}{
  \usoquando{Si vuole il valore «attorno a cui» si distribuisce la variabile: il valore atteso, il baricentro della
    pmf.}
  \usocome{Si moltiplica ogni valore per la sua probabilità e si somma.}
  \usoattenzione{La media non è per forza un valore che la variabile può assumere: nessun dado mostra $3{,}5$.}
  \usoesempio{Dado onesto: $E[X] = \frac{1 + 2 + \cdots + 6}{6} = 3{,}5$. Dado truccato con $P(1) = \frac{5}{6}$ e
    $\frac{1}{30}$ alle altre facce: $E[X] = \frac{5}{6} + \frac{2 + 3 + 4 + 5 + 6}{30} = 1{,}5$.}
}

\scheda{bernoulli}{
  \usoquando{Una sola prova con due esiti: un bit giusto o sbagliato, testa o croce, un pezzo buono o difettoso.}
  \usocome{Si chiama «successo» (valore 1) l'esito che interessa, e $p$ la sua probabilità.}
  \usoesempio{Un bit è sbagliato con probabilità $0{,}1$: la variabile «il bit è sbagliato» ha $E[X] = 0{,}1$ e
    $\sigma_X^2 = 0{,}1 \cdot 0{,}9 = 0{,}09$.}
}

\scheda{distbinomiale}{
  \usoquando{Si contano i successi in $N$ prove \emph{indipendenti}, tutte con la \emph{stessa} probabilità di
    successo $p$: bit sbagliati in un pacchetto, teste in $N$ lanci.}
  \usocome{Si individuano $N$, $p$ e il numero $k$ di successi richiesto: $\binom{N}{k}$ conta le sequenze con $k$
    successi, $p^k (1-p)^{N-k}$ è la probabilità di ciascuna.}
  \usovarianti{Per «almeno un successo» si passa al complementare: $1 - (1 - p)^N$. Per «almeno $k$» si sommano i
    casi da $k$ a $N$, oppure si toglie da 1 la somma dei casi da 0 a $k - 1$: conviene la somma più corta.}
  \usoattenzione{Servono prove indipendenti e $p$ costante. Le estrazioni \emph{senza} reinserimento non sono
    binomiali: lì si contano i casi con le combinazioni (\vv{combinazioni}).}
  \usoesempio{Pacchetto di 10 bit, ognuno sbagliato con probabilità $0{,}1$:
    $P(X = 2) = \binom{10}{2}\, 0{,}1^2 \cdot 0{,}9^8 \simeq 0{,}1937$ e $P(X \geq 1) = 1 - 0{,}9^{10} \simeq
    0{,}6513$. In media $E[X] = 1$ bit sbagliato, con $\sigma_X^2 = 0{,}9$.}
}

\scheda{uniforme}{
  \usoquando{La variabile assume $M$ valori, tutti con la stessa probabilità: dado onesto, cifra casuale, byte
    casuale.}
  \usocome{La media è la media aritmetica dei valori; su $\{0, \ldots, M - 1\}$ ci sono le formule pronte.}
  \usoesempio{Un byte casuale ($M = 256$): $E[X] = 255/2 = 127{,}5$ e $\sigma_X^2 = (256^2 - 1)/12 = 5\,461{,}25$.}
}

\scheda{poisson}{
  \usoquando{Si contano eventi rari in un intervallo di tempo o di spazio, e se ne conosce il numero medio
    $\lambda$: pacchetti che arrivano a un router in un millisecondo, guasti in un mese. Oppure si vuole approssimare
    una binomiale con $N$ grande e $p$ piccolo.}
  \usocome{Si fissa $\lambda$, il numero medio di eventi \emph{nell'intervallo considerato}, e si applica la pmf.
    Per approssimare la binomiale si prende $\lambda = N p$.}
  \usoattenzione{Se cambia l'intervallo cambia anche $\lambda$: con 3 arrivi al millisecondo in media, in due
    millisecondi la media è 6.}
  \usoesempio{$\lambda = 3$: $P(X = 0) = e^{-3} \simeq 0{,}0498$, quindi $P(X \geq 1) \simeq 0{,}9502$;
    $P(X = 3) = \frac{3^3}{3!}\, e^{-3} \simeq 0{,}2240$. La binomiale $\mathcal{B}(30;\ 0{,}1)$, con la stessa
    media, dà $P(X = 3) \simeq 0{,}2361$.}
}
